% Generated draft: requires mathematical and source review.
\documentclass[11pt]{article}
\usepackage[margin=25mm]{geometry}
\usepackage{fontspec}
\setmainfont{Latin Modern Roman}
\usepackage{amsmath,amssymb,mathtools,unicode-math}
\setmathfont{Latin Modern Math}
\usepackage{enumitem,needspace,xurl,hyperref}
\setlength{\parindent}{0pt}
\setlength{\parskip}{5pt}
\setlength{\emergencystretch}{3em}
\newcommand{\N}{\mathbb{N}}
\newcommand{\Z}{\mathbb{Z}}
\newcommand{\Q}{\mathbb{Q}}
\newcommand{\R}{\mathbb{R}}
\newcommand{\C}{\mathbb{C}}
\newcommand{\F}{\mathbb{F}}
\newcommand{\A}{\mathbb{A}}
\newcommand{\PP}{\mathbb{P}}
\newcommand{\E}{\mathbb{E}}
\newcommand{\eps}{\varepsilon}
\newcommand{\dd}{\,\mathrm{d}}
\begin{document}
\section*{UnsolvedMath ID: 30003199}
\textbf{Source problem code:} catalogue entry
\subsection{Definitions and mathematical statement}
Let \(\mathcal Z\) denote the Jiang--Su algebra. A \(C^*\)-algebra \(A\) is separable if it has a countable norm-dense subset, unital if it has an identity \(1_A\), and simple if it is nonzero and has no nonzero proper closed two-sided ideals. It is nuclear if, for every \(C^*\)-algebra \(D\), the minimal and maximal \(C^*\)-tensor norms on \(A\odot D\) coincide. It is \(\mathcal Z\)-stable if \(A\cong A\otimes_{\min}\mathcal Z\). Define
\[
\mathscr C=\{A: A\text{ is a nonzero separable, unital, simple, nuclear }C^*\text{-algebra and }A\not\cong A\otimes_{\min}\mathcal Z\}.
\]
The original wording uses ``amenable'' where the cleaned statement uses ``nuclear''; these are equivalent properties for \(C^*\)-algebras. No Universal Coefficient Theorem hypothesis or additional regularity hypothesis is supplied.

For a unital \(C^*\)-algebra \(A\), write the usual Elliott invariant as
\[
\operatorname{Ell}(A)=\bigl(K_0(A),K_0(A)^+,[1_A],K_1(A),T(A),r_A\bigr).
\]
Here \(K_0(A)\) is the Grothendieck group formed from stable Murray--von Neumann equivalence classes of projections, \(K_0(A)^+\) is its projection cone, \([1_A]\) is the unit class, and \(K_1(A)\) is the group of stable homotopy classes of unitaries. The space \(T(A)\) consists of tracial states with its weak-* topology and convex structure; it may be empty. The pairing is
\[
r_A(\tau,[p]-[q])=(\tau\otimes\operatorname{Tr})(p)-(\tau\otimes\operatorname{Tr})(q),
\]
where \(\operatorname{Tr}\) is the unnormalized matrix trace. An isomorphism of these invariants preserves the ordered \(K_0\)-group and unit, the \(K_1\)-group, the topological affine structure of the trace spaces, and this pairing.

Let \(\mathcal K\) be the compact operators on a separable infinite-dimensional Hilbert space. For \(a,b\in(A\otimes\mathcal K)_+\), write \(a\precsim b\) if there are \(x_n\in A\otimes\mathcal K\) such that \(\|x_n^*bx_n-a\|\to0\). Write \(a\sim b\) if both \(a\precsim b\) and \(b\precsim a\). The stabilized Cuntz semigroup is
\[
\operatorname{Cu}(A)=(A\otimes\mathcal K)_+/\!\sim,
\]
with order induced by \(\precsim\) and addition induced by block direct sum. Its standard structure includes suprema of increasing sequences and the relation \(x\ll y\): whenever \(y\leq\sup_n z_n\) for an increasing sequence, \(x\leq z_n\) for some \(n\). Isomorphisms of Cuntz semigroups preserve addition, order, increasing sequential suprema, and this relation.

The first source question asks whether the Elliott invariant together with the Cuntz semigroup classifies the algebras in \(\mathscr C\). At the level of determining isomorphism classes, a literal interpretation using the two invariants separately asks whether, for every \(A,B\in\mathscr C\),
\[
\bigl(\operatorname{Ell}(A)\cong\operatorname{Ell}(B)\ \text{and}\ \operatorname{Cu}(A)\cong\operatorname{Cu}(B)\bigr)
\quad\Longrightarrow\quad A\cong B.
\]
The reverse implication follows from invariance. This displayed implication specifies the bare-pair interpretation only. The source does not specify whether ``together'' requires compatibility between the two invariant isomorphisms, distinguished Cuntz-semigroup data, or additional structure linking projections and traces to the Cuntz semigroup. It also does not specify whether classification must lift every appropriate invariant isomorphism to an algebra isomorphism. Those stronger or differently structured interpretations remain unresolved formulation choices.

The second source question asks what classification is possible in the nonsimple case. Dropping simplicity from the preceding hypotheses gives one possible comparison class, but the record does not explicitly settle whether non-\(\mathcal Z\)-stability is retained, which nonsimple subclasses are intended, or which invariant and notion of classification should be used. This part is therefore retained as an open-ended classification question rather than assigned an unsupported universal assertion.
\subsection{Short English statement}
For separable, unital, simple, nuclear \(C^*\)-algebras that are not \(\mathcal Z\)-stable, ask whether the Elliott invariant together with the Cuntz semigroup determines the algebra up to isomorphism. Also ask what classification is possible without simplicity. The intended compatibility of the invariants and the precise nonsimple scope are unspecified. The supplied record labels the problem partially solved; this draft does not independently verify that assessment.
\subsection{Sources}
\begin{itemize}
\item Local UnsolvedMath registry: supplied complete record for permanent ID \(30003199\), source problem code \texttt{catalogue entry}, in the Oberwolfach Reports Open Problems collection. No individual proposer is supplied; the proposed year is recorded as 2016.
\item Supplied report citation: \(C^*\)-Algebras (2017), \emph{Oberwolfach Reports}, pp.\ 2269--2345, DOI \texttt{10.4171/owr/2016/40}. \url{https://doi.org/10.4171/owr/2016/40}. The record describes these reports as documenting mathematical workshops and open-problem discussions. Its citation uses 2017, while its proposed year and year tag use 2016; this discrepancy is not resolved here.
\item Supplied literature source, labelled ``Cantier Unitary Cuntz semigroup'': \url{https://arxiv.org/abs/2111.06926}.
\item Supplied literature source, labelled ``Gardella and Perera Cuntz semigroups survey'': \url{https://doi.org/10.4171/EMSS/84}.
\item Supplied literature source, labelled ``OWR Operator Algebras report'': \url{https://ems.press/content/serial-article-files/46645?nt=1}.
\item Literature assessment dated exactly \texttt{2026-08-21}: the supplied curation labels the status ``partially solved'' and states that Elliott and Cuntz semigroup methods classify substantial simple or regular subclasses, while classification in the stated non-\(\mathcal Z\)-stable and nonsimple generality remains unresolved. This is supplied dated curation derived from the OWR material, not an independently verified literature assessment; the record expressly calls for independent verification of time-sensitive status and cited claims.
\item The supplied statement-status label is \texttt{corrected\_verified}. Its verification note reports that the official report was checked and that the clean statement preserves its classification hypotheses and two follow-up questions. This is a claim of the supplied registry, not a source visit or verification performed for this draft.
\end{itemize}
\end{document}
